\section{Proposed Placement and Admission}
We provisionally call the mechanism Service-Aware Placement and Admission
(\mechanism). The initial implementation would combine an offline
placement planner with finite endpoint transaction budgets, keeping
packet packing and VC policies fixed.

\subsection{Planner Inputs and Decisions}
Inputs include object sizes, observed access counts, stage dependencies,
per-die free HBM, Host and MPU capacity, calibrated service curves, and
the mapping of traffic onto shared paths. Decisions assign objects or
shards to memory tiers and choose endpoint outstanding windows for
traffic classes. The first version avoids replication to keep capacity
accounting and consistency manageable.

A practical greedy planner can evaluate candidate migrations by the HBM
space released and their incremental demand on each constrained resource.
It prioritizes feasible migrations by released capacity per unit of
dominant-resource pressure, handling zero incremental demand separately.
After each assignment it recomputes shared-path demand. Hot objects are
sharded only when the resulting service parallelism outweighs additional
transactions and path overhead. The objective is to enlarge the admitted
working set under the budget, rather than obtain low latency merely by
serving fewer requests. Small-instance exhaustive search provides a
reference for the quality of this heuristic.

\subsection{Transaction Admission Boundary}
The endpoint acquires an outstanding-transaction token before issuing
remote work. Response storage must be finite and accounted for; the
proposed token is released after the associated data reaches the chosen
HBM-consumption point, as illustrated in Fig.~\ref{fig:lifetime}.
Internal MPU slots remain separate and can be released earlier.
Admission tokens are not CBFC credits and do not change protocol credit
semantics, return-port mapping, or transaction ordering.

The implementation must specify response reservations, bounded waiting,
and resource release on completion or failure. New requests must not
consume all resources needed by responses. Progress tests are necessary;
a smaller window alone is not a deadlock argument. Planner overhead,
state size, and sensitivity to inaccurate access estimates belong in
the evaluation.

\subsection{Scope of the First Version}
Service curves are fitted using calibration workloads, then frozen
before evaluating unseen traces. A dynamic controller is unnecessary
until static configurations establish useful headroom. If a workload
cannot satisfy the budget, the policy may choose a different tier or
reject additional work; both rejection and unserved demand are visible
outcomes. Whole-system service guarantees require an arrival and batching
model beyond the current representative-stage simulator.
